% Generated draft: requires mathematical and source review.
\documentclass[11pt]{article}
\usepackage[margin=25mm]{geometry}
\usepackage{fontspec}
\setmainfont{Latin Modern Roman}
\usepackage{amsmath,amssymb,mathtools,unicode-math}
\setmathfont{Latin Modern Math}
\usepackage{enumitem,needspace,xurl,hyperref}
\setlength{\parindent}{0pt}
\setlength{\parskip}{5pt}
\setlength{\emergencystretch}{3em}
\newcommand{\N}{\mathbb{N}}
\newcommand{\Z}{\mathbb{Z}}
\newcommand{\Q}{\mathbb{Q}}
\newcommand{\R}{\mathbb{R}}
\newcommand{\C}{\mathbb{C}}
\newcommand{\F}{\mathbb{F}}
\newcommand{\A}{\mathbb{A}}
\newcommand{\PP}{\mathbb{P}}
\newcommand{\E}{\mathbb{E}}
\newcommand{\eps}{\varepsilon}
\newcommand{\dd}{\,\mathrm{d}}
\begin{document}
\section*{UnsolvedMath ID: 30004657}
\textbf{Source problem code:} catalogue entry
\subsection{Definitions and mathematical statement}
The supplied record poses three construction questions rather than a fully specified conjecture. The following standard terminology makes their intended content explicit without supplying missing hypotheses.

In the usual smooth, finite-dimensional setting, let \(Q\) be a configuration manifold and let \(L:TQ\to\R\) be a Lagrangian. In local coordinates \((q^i,v^i)\), its velocity Hessian is
\[
W_{ij}(q,v)=\frac{\partial^2 L}{\partial v^i\partial v^j}(q,v).
\]
Degeneracy means singularity of this Hessian on the region under consideration. The source does not specify that region, a constant-rank condition, or hypotheses ensuring consistent constraints and a well-defined evolution. This terminology therefore does not fix the intended class of general degenerate Lagrangians.

A symplectic potential is a one-form \(\theta\) for a symplectic form \(\omega\); using the convention \(\omega=-\dd\theta\), write locally \(\theta=\sum_i\theta_i(z)\,\dd z^i\). A velocity-linear Lagrangian associated with such a potential has the form
\[
L(z,\dot z)=\theta_z(\dot z)-H(z).
\]
This is an illustrative class, not an asserted representation of every degenerate Lagrangian. The existence of a local or global potential and its relation to the general Lagrangians in the question are not specified by the source. The condition that a component \(\theta_i\) vanish depends on the chosen coordinates and potential.

For terminology, a discrete variational principle seeks stationary points of a discrete action. In the basic endpoint formulation, a discrete Lagrangian \(L_d\) gives
\[
S_d(q_0,\ldots,q_N)=\sum_{k=0}^{N-1}L_d(q_k,q_{k+1};h),
\qquad \delta S_d=0,
\qquad \delta q_0=\delta q_N=0.
\]
Galerkin and Runge--Kutta variational constructions may additionally involve internal stages and constraints. This displayed endpoint formulation is only a definition of standard terminology; the source does not require it as the particular principle sought. A numerical map \(\Phi_h\) is symplectic on its specified phase space if \(\Phi_h^*\omega=\omega\). When a smooth exact evolution \(\varphi_h\) exists, order at least \(p\) conventionally means local error \(\Phi_h(z)-\varphi_h(z)=O(h^{p+1})\) in local coordinates under suitable regularity assumptions. For constrained dynamics, the relevant phase space and admissible initial data must first be specified.

\textbf{Questions recorded in the source.}
\begin{enumerate}
\item For general degenerate Lagrangians, can degenerate variational integrators be extended to arbitrary order by Galerkin or Runge--Kutta discretizations? Here arbitrary order asks for a construction for every positive integer \(p\). The required discrete degeneracy property and admissible class of continuous systems remain unspecified.
\item Can these degenerate variational integrators and symplectic Runge--Kutta methods be obtained without imposing that any component of the symplectic potential vanish? The intended coordinates, potential, phase space, and sense of obtaining the methods must be specified before this becomes a precise construction claim.
\item Do symplectic Runge--Kutta methods follow from a discrete variational principle? The record does not specify the quantifier over methods, the applicable class of systems, the admissible form of the principle, or the precise equivalence required between its stationary equations and the numerical method.
\end{enumerate}
These questions retain the source's generality. No universal existence or equivalence theorem is asserted, and the missing qualifications prevent a unique fully quantified formulation.
\subsection{Short English statement}
The problem asks for arbitrary-order variational integrators for general degenerate Lagrangians, constructions that impose no vanishing component of a symplectic potential, and a discrete variational derivation of symplectic Runge--Kutta methods. The source leaves the system class, constraints, discrete degeneracy condition, and intended equivalence unspecified. Its open status is supplied dated curation, not an independently verified conclusion.
\subsection{Sources}
\begin{itemize}
\item Supplied local UnsolvedMath registry record, permanent ID \(30004657\), title \emph{High-Order Variational Integrators for Degenerate Lagrangians}, exact source problem code \texttt{catalogue entry}, in the Oberwolfach Reports Open Problems collection. The supplied proposed year is 2021; no individual proposer is named.
\item Supplied report citation: \emph{Geometric Numerical Integration (2022)}, Oberwolfach Reports, DOI 10.4171/owr/2021/17, pp. 943--999. \url{https://doi.org/10.4171/owr/2021/17}. The background describes Oberwolfach Reports as documenting mathematical workshops and their open-problem discussions. The supplied title year 2022 and report designation 17/2021 are retained without resolving their relationship.
\item Supplied literature reference: Leok--Ohsawa--Sosa, variational partitioned Runge--Kutta methods. \url{https://arxiv.org/abs/1401.7904}.
\item Supplied literature reference: Sanz-Serna, discrete variational methods. \url{https://arxiv.org/abs/2001.07185}.
\item Supplied literature reference: Oberwolfach Report 17/2021. \url{https://doi.org/10.4171/owr/2021/17}.
\item Literature assessment dated exactly \texttt{2026-08-22}: the supplied curation labels the problem open and reports existing variational and symplectic constructions for important degenerate or velocity-linear classes, while stating that the requested general constructions and equivalence result are not established in the stated generality. This is supplied dated curation derived from the OWR material, not independently verified literature assessment. The source explicitly requests independent verification of time-sensitive status and cited claims.
\item The registry labels the statement \texttt{corrected\_verified}, explaining that a publication-format audit preserved mathematical content while repairing notation or typography. This supplied audit description does not establish mathematical correctness or resolve the formulation gaps above.
\end{itemize}
\end{document}
